% !TEX root = ../main.tex
\corrigesChapitre{2}{Suites numériques}
\corrigesSouples

\begin{blocLivre}{cor-2-1}
\begin{corrige}{1}{Substituer un indice}
\demarche{Remplacer chaque occurrence de l’indice par la même valeur.}
\begin{questions}
\item On remplace l’indice dans les trois termes : $u_0=0-0+2=\res{2}$,
$u_1=1-3+2=\res{0}$, $u_2=4-6+2=\res{0}$ et $u_5=25-15+2=\res{12}$.
\item $u_{n+1}=(n+1)^2-3(n+1)+2=n^2+2n+1-3n-3+2=n^2-n$.
On soustrait ensuite l’expression entière de $u_n$ :
\[u_{n+1}-u_n=(n^2-n)-(n^2-3n+2)=\res{2n-2}.\]
Les parenthèses évitent de perdre le signe moins devant $u_n$.
\end{questions}

\end{corrige}
\end{blocLivre}
\begin{blocLivre}{cor-2-2}
\begin{corrige}{2}{Calculer de proche en proche}
\demarche{Calculer dans l’ordre : chaque terme réutilise le terme précédent.}
\begin{questions}\item $v_1=3\times2-1=5$, puis $\res{v_2=14,\ v_3=41,\ v_4=122}$.
\item Les différences sont $3$, $9$, $27$ ; les rapports sont $5/2$, $14/5$, $41/14$.
Ils ne sont pas constants : la suite n'est ni arithmétique ni géométrique.
\end{questions}
\end{corrige}
\end{blocLivre}
\begin{blocLivre}{cor-2-3}
\begin{corrige}{3}{Une suite arithmétique}
\demarche{Employer la formule arithmétique, puis résoudre une inégalité sur les indices entiers.}
\begin{questions}\item $\res{u_n=7-3n}$ et $u_8=-17$. Comme la raison est négative, la suite est strictement décroissante.
\item $7-3n\leq-20\iff n\geq9$ : les indices cherchés sont tous les entiers $\res{n\geq9}$.
\end{questions}
\end{corrige}
\end{blocLivre}
\begin{blocLivre}{cor-2-4}
\begin{corrige}{4}{Une suite géométrique}
\demarche{Employer la formule géométrique et compter les termes de la somme.}
\begin{questions}\item $\res{v_n=5\times2^n}$, $v_4=80$ et $v_6=320$.
\item Directement : $5+10+20+40+80=155$.
Avec la formule : $5(1-2^5)/(1-2)=\res{155}$. Il y a cinq termes.
\end{questions}
\end{corrige}
\end{blocLivre}
\begin{blocLivre}{cor-2-5}
\begin{corrige}{5}{Du pourcentage à la formule}
\demarche{Transformer la baisse annuelle en coefficient multiplicateur.}
\begin{questions}\item Le coefficient multiplicateur est $1-0{,}15=0{,}85$ :
$V_{n+1}=0{,}85V_n$, $V_0=800$, donc $\res{V_n=800\times0{,}85^n}$.
\item $V_2=578$ euros. La baisse vaut $1-0{,}85^2=0{,}2775$, soit
$\res{27{,}75\,\%}$. Les pourcentages successifs se multiplient via leurs coefficients.
\end{questions}
\end{corrige}
\end{blocLivre}
\begin{blocLivre}{cor-2-6}
\begin{corrige}{6}{Retrouver les paramètres}
\demarche{Compter les transitions entre les deux indices connus.}
\begin{questions}
\item \methode{utiliser l'écart entre les indices}
$u_8-u_3=5r=20$, donc $r=4$ ; $u_0=u_3-3r=-1$.
Ainsi $\res{u_n=4n-1}$.
\item $v_5=v_2q^3$ donne $q^3=8$, donc $q=2$.
$v_0=12/2^2=3$, d'où $\res{v_n=3\times2^n}$.
\end{questions}
\end{corrige}
\end{blocLivre}
\begin{blocLivre}{cor-2-7}
\begin{corrige}{7}{Compter les termes d'une somme}
\demarche{Identifier premier terme, dernier terme et nombre de termes avant de sommer.}
\begin{questions}
\item $61=4+19\times3$ : vingt termes, de rang 0 à 19.
La somme vaut $20(4+61)/2=\res{650}$.
\item $384=3\times2^7$ : huit termes. La somme vaut
$3(2^8-1)=\res{765}$.
\item De 3 à 10 inclus, huit termes ; le premier vaut 7, le dernier 21.
La somme est $8(7+21)/2=\res{112}$.
\end{questions}
\end{corrige}
\end{blocLivre}
\begin{blocLivre}{cor-2-8}
\begin{corrige}{8}{Comparer des variations}
\demarche{Calculer la différence de deux termes consécutifs pour étudier les variations.}
\begin{questions}
\item $u_{n+1}-u_n=2n-3$, négatif pour $n=0,1$, positif pour $n\geq2$.
La suite décroît jusqu'au rang 2, puis croît ; elle n'est pas monotone sur tous les indices.
\item Au même dénominateur,
$v_{n+1}-v_n=\dfrac{(n+2)^2-(n+1)(n+3)}{(n+2)(n+3)}
=\dfrac1{(n+2)(n+3)}>0$. La suite est strictement croissante.
\item $v_n=(n+2-1)/(n+2)=1-1/(n+2)$.
Les termes restent inférieurs à 1 et s'en rapprochent quand $n$ devient grand.
Il s'agit ici d'une lecture intuitive du comportement à long terme.
\end{questions}
\end{corrige}
\end{blocLivre}
\begin{blocLivre}{cor-2-9}
\begin{corrige}{9}{Lire puis justifier un nuage de points}
\demarche{Séparer la conjecture graphique de la justification donnée par la récurrence.}
\begin{questions}
\item Les termes semblent divisés par 2 à chaque étape : $u_n=16(1/2)^n$.
\item La récurrence donnée est précisément celle d'une suite géométrique
de raison $1/2$, de premier terme 16, ce qui justifie la formule.
Les termes sont positifs et la raison est dans $]0;1[$ : décroissance stricte.
\item $u_7=0{,}125$ et $u_8=0{,}0625$. Par décroissance, le premier rang
tel que $u_n<0{,}1$ est $\res{n=8}$.
\end{questions}
\end{corrige}
\end{blocLivre}
\begin{blocLivre}{cor-2-10}
\begin{corrige}{10}{Comprendre une boucle Python}
\demarche{Distinguer le nombre de transitions et le nombre de valeurs stockées.}
\begin{questions}
\item Les valeurs successives sont 3, 7, 15, 31, 63. La liste finale est
\code{[3, 7, 15, 31, 63]}.
\item $u_0=3$, $u_{n+1}=2u_n+1$. Après quatre passages, on a calculé $u_4$.
\item Remplacer \code{range(4)} par \code{range(10)}.
La liste comprend alors onze termes, de $u_0$ à $u_{10}$.
\end{questions}
\end{corrige}
\end{blocLivre}
\begin{blocLivre}{cor-2-11}
\begin{corrige}{11}{Un seuil d'épargne}
\demarche{Justifier la croissance, puis vérifier le seuil et son prédécesseur.}
\begin{questions}
\item La raison vaut $1{,}05>1$ et le premier terme est positif : croissance stricte.
\item On initialise $n=0$ et $C=1000$, puis on actualise ces deux variables ensemble.
\end{questions}
\begin{python}
n = 0
C = 1000
while C < 1500:
    C = 1.05*C
    n = n + 1
\end{python}
\otLabel{3.}\enspace $C_8\approx1477{,}46$ et $C_9\approx1551{,}33$ euros.
La première année cherchée est donc $\res{n=9}$.

À chaque test, les variables représentent bien $C_n$ et son rang $n$. L’arrêt à 9 donne un candidat ; la croissance et la vérification de $C_8$ prouvent que ce candidat est le premier.

\end{corrige}
\end{blocLivre}
\begin{blocLivre}{cor-2-12}
\begin{corrige}{12}{Des rangées de sièges}
\demarche{Respecter le premier indice 1 et additionner des rangées consécutives.}
\begin{questions}
\item \methode{ici le premier indice est 1}
$u_n=18+2(n-1)=\res{2n+16}$, donc $u_{20}=56$.
\item La capacité est $20(18+56)/2=\res{740\text{ places}}$.
\item Les dix premières rangées contiennent $10(18+36)/2=270$ sièges ;
les dix dernières, $10(38+56)/2=470$.
La recette vaut $9\times270+12\times470=\res{8070\text{ euros}}$.
\item Avec $N$ rangées, la capacité est $S_N=N(18+2N+16)/2=N(N+17)$.
Elle croît avec $N\geq1$. Or $S_{24}=984<1000$ et $S_{25}=1050\geq1000$.
Il faut donc au moins $\res{25\text{ rangées}}$.
\end{questions}
\end{corrige}
\end{blocLivre}
\begin{blocLivre}{cor-2-13}
\begin{corrige}{13}{Un traitement répété}
\demarche{Soustraire la valeur d’équilibre pour obtenir une suite géométrique.}
\begin{questions}
\item On conserve $0{,}8u_n$, puis on ajoute 6 : $u_{n+1}=0{,}8u_n+6$.
$u_1=14$ mg et $u_2=17{,}2$ mg.
\item $v_{n+1}=u_{n+1}-30=0{,}8u_n-24=0{,}8(u_n-30)=0{,}8v_n$.
\item $v_0=-20$, donc $v_n=-20\times0{,}8^n$ et
$\res{u_n=30-20\times0{,}8^n}$.
Puis $u_{n+1}-u_n=4\times0{,}8^n>0$ et $30-u_n=20\times0{,}8^n>0$.
\item $u_{10}\approx27{,}853<28$ et $u_{11}\approx28{,}282\geq28$.
Par croissance, le premier indice est $\res{n=11}$.
\end{questions}
\label{figlivre-r02-13-cycle-cor}
\figureLivre[corrige]{r02-13-cycle}

\end{corrige}
\end{blocLivre}
\begin{blocLivre}{cor-2-14}
\begin{corrige}{14}{Deux contrats de rémunération}
\demarche{Distinguer indice de la suite et numéro de l’année.}
\begin{questions}
\item $A_n=24000+1000n$ et $B_n=24000\times1{,}04^n$.
La sixième année correspond à $n=5$ : $A_5=29000$ euros et
$B_5\approx29199{,}67$ euros ; B est supérieur.
\item Les dix premières années vont de $n=0$ à $n=9$.
\[\sum_{n=0}^9 A_n=10\frac{24000+33000}2=\res{285000},\]
\[\sum_{n=0}^9 B_n=24000\frac{1{,}04^{10}-1}{0{,}04}\approx\res{288146{,}57}.\]
B rapporte environ 3146,57 euros de plus sur dix ans.
\item Pour $n=1,2,3$, les salaires B valent respectivement
24960 ; 25958,40 ; 26996,736 euros, inférieurs à 25000 ; 26000 ; 27000.
Pour $n=4$, $B_4=28076{,}60544>28000=A_4$.
Le premier indice est donc $\res{n=4}$, soit la cinquième année.
\end{questions}
\label{figlivre-r02-14-salaires-cor}
\figureLivre[corrige]{r02-14-salaires}

\end{corrige}
\end{blocLivre}
\begin{blocLivre}{cor-2-15}
\begin{corrige}{15}{Des carrés construits avec des allumettes}
\demarche{Compter les allumettes nouvelles, puis additionner les constructions séparées.}
\begin{questions}
\item Le premier carré utilise 4 allumettes ; chaque carré supplémentaire en ajoute 3.
Ainsi $a_1=4$, $a_{n+1}=a_n+3$ et $\res{a_n=3n+1}$.
\item $3n+1\leq100\iff n\leq33$ : on peut aligner $\res{33\text{ carrés}}$,
avec exactement 100 allumettes.
\item Les constructions sont séparées : $T_N=\sum_{n=1}^N(3n+1)$.
Donc $T_N=3N(N+1)/2+N=\res{(3N^2+5N)/2}$.
\item $T_{25}=1000$ et $T_{26}=1079$. Comme $T_N$ est croissant,
le maximum possible est $\res{N=25}$.
\end{questions}
\end{corrige}
\end{blocLivre}
